\scanpage{125}

\uline{Dévissage de $\mathbb Q_\ell$-faisceaux}

\deleted{Un faisceau}

Un $\mathbb Q_\ell$-faisceau \deleted{\unsure{est dit}} \inserted{$F$} sur $X$ est dit \uline{admissible} \inserted{de type $\leq d$} s'il satisfait la condition suivante :

$\forall x\in X$, existe voisinage ouvert $U$ de $x$ dans $\overline{x}$ et filtration \inserted{finie} de $F|U$ par des ss-faisceaux \deleted{\unsure{$F^i$}} tels que les $\operatorname{gr}^i(F|U)$ soient isomorphes à un \deleted{``facteur''} \inserted{quotient \inserted{\unsure{ess. constant|ens. constant|en. constituants}} d'un sous-objet ess. constant} de un faisceau de la forme $R^n g_*(\deleted{\unsure{\mathbb Q_{\ell V}}}\mathbb Q_\ell)$, où $g:V\unsure{\longrightarrow|\text{---}}U$ est un morphisme \uline{projectif}, \deleted{\unsure{et}} \uline{lisse}, à fibres connexes \inserted{de dim $\leq d$}.

% slanted
\margin{``ess. constant'' signifie : faisceau de \unsure{Steenrodien|Steenro-dien} \unsure{!|''}}

[Se vérifie fibre par fibre, + $F$ est ``constructible'']

\deleted{Ceci dit que $F$ est \unsure{admis. si|admissible}} \inserted{\unsure{vrai} pour \uline{fibres} \uline{géom}}

\uline{Lemme} 1 \quad Les \deleted{\unsure{fai}} $\mathbb Q_\ell$-faisceaux admissibles de type $\leq d$ forment une sous-catégorie abélienne \deleted{\unsure{exacte} et épaisse (stable par} \inserted{épaisse de la catégorie des $\mathbb Q_\ell$-faisceaux constructibles} \deleted{noyau, conoyau, extension) de la catégorie des}

\deleted{\unsure{2}} (i) Stable par ss-\deleted{\unsure{fais|faisc.}}objet et objet quotient.

En effet, si $F|U$ est muni d'une filtration $d$-spéciale, \inserted{et si $G|U$ est ess. constant,} alors la filtration induite sur $G$ resp $(F\unsure{\vert U|\cap U})/G$ est $d$-spéciale. \deleted{\unsure{car $H$ de}}

(ii) Stable par extensions : trivial.

\uline{Lemme} 2 \quad \uline{Images inverses} \qquad trivial.

% archive sheet number: 124
